%======================================================================
%  Results -- Model interpretation  (paper section 5; main figure = Figure 4)
%
%  Source artifact: GPN-Plasmids Interpretation
%    https://claude.ai/code/artifact/fa19d880-b92e-4e9f-94d0-6743e313a07a
%    /home/canall/DNA_LM/plasmids_GPN/manuscript/artifact_interpretation/page.html
%    snapshot 2026-08-28 (page.html mtime 2026-08-28 15:47)
%
%  Eight paragraphs, no subsection headings, from the artifact's "Main text"
%  band. The author has frozen this section's wording ("come back later"), so
%  it is transcribed exactly -- including passages that read oddly.
%
%  The method-writeups/*.md files in that tree are STALE and contradict the
%  HTML; they were not used. This artifact also ships no bibliography and no
%  citation markers. See CONVERSION_NOTES.md.
%
%  NOTE: the artifact's section id is still "figure5" and its assets are all
%  named fig5_*, but its label is already "Figure 4"; the paper numbers it
%  Figure 4. Of this artifact's Supplementary S1-S7, S5 was promoted to
%  Extended Data Figure 5 (2026-09-01); the rest renumber globally to S9-S14.
%======================================================================

\subsection*{Model interpretation}
\label{sec:interpretation}

Having applied the \model classifiers across IMG/PR
(\secref{sec:genomewide}), we examined the frozen pretrained model directly,
% TODO wording: the artifact reads "across IMG/PR (§4)".
to ask which recurring sequence dependencies it had acquired during
pretraining. We measured these with a categorical Jacobian: within a fixed
2,048 bp context, each position in turn was substituted with all four
nucleotides and the resulting change in the model's predicted logits was
recorded at every other position, giving a finite-difference measure of how
far the prediction at one site depends on the identity of another
(\nameref{sec:methods}). The construction is the one used to extract
co-evolutionary signal from protein and genomic language models,
% TODO cite: the categorical-Jacobian construction (protein/genomic LM
% co-evolution). Named in prose only; the artifact ships no bibliography.
and we refer to its output as a categorical-Jacobian dependency map. It is a
property of the frozen model's predictive distribution: it is not a
measurement of physical contact, base pairing, or evolutionary covariation,
and no annotation enters the computation.

Two visual signatures recur across the loci examined. Direct repeats are
associated with off-diagonal bands running parallel to the main diagonal at
the repeat period, and inverted repeats with anti-diagonal patterns centred
on the paired arms. We report these as recurring in the examples shown rather
than as a demonstrated class distinction: the loci here were selected because
their architecture is documented, no quantitative test was run across a
defined denominator of loci, and every panel carries its own local colour
scale, so pattern intensities cannot be compared between panels.

\fref{fig:interpretation} shows two characterized case studies. At the
minimal replication origin of \emph{Escherichia coli} plasmid pSC101
(NC\_002056.1, displayed crop [5100,5272)), the three 21 bp RepA iterons
DR1--DR3 are associated with a parallel off-diagonal ladder in which the
DR1$\times$DR2 and DR1$\times$DR3 bands are separately resolved, while the
inverted repeats IR1 and IR2 show anti-diagonal patterns; the 11 bp spacer
between the second and third iterons, reported to tolerate substitution
without loss of origin function, is comparatively featureless
(\fref{fig:interpretation}b). The displayed interval is a crop within the
reported minimal origin: the DnaA box and IHF site lie outside it and are not
shown. At the origin of transfer of \emph{Klebsiella pneumoniae} plasmid
pIGRK (AY543071.1, [904,1214)), IR1, IR2 and IR3 each show a discrete
anti-diagonal pattern and the 12 bp direct repeat DR1 a compact on-diagonal
block, against a comparatively featureless background; signal also occurs
across the P\emph{mobK} $-$35/$-$10 box pair, which the source describes as a
promoter predicted computationally rather than one mapped experimentally
(\fref{fig:interpretation}a). In both panels the model-derived dependency
patterns align with curated repeat and promoter annotations. Two loci do not
establish that this holds for origins generally.

Supplementary panels extend the observation to further locus classes, all of
them curated or hand-selected examples rather than an unbiased sample. Eight
additional origins --- four of transfer and four of replication, spanning
three architectural classes --- show the same association between repeat
geometry and dependency-pattern geometry
(\suppfrefs{fig:supp-orit-maps}{fig:supp-oriv-maps}). Across four
$\sigma^{70}$ promoters, off-diagonal signal is present between the $-$35 and
$-$10 boxes across spacings of 17--18 bp, the Shine--Dalgarno region carries
a dense on-diagonal block, and the coding sequence shows periodic on-diagonal
beading at codon spacing (\suppfref{fig:supp-promoter}). These four were
chosen as visually clean examples from a motif scan and do not support a
promoter-wide recovery rate. Long periodic direct-repeat arrays behave
comparably: in four CRISPR arrays the direct repeats are associated with a
periodic off-diagonal lattice at the repeat spacing
(\suppfref{fig:supp-crispr}). We include the CRISPR panels as evidence that
the direct-repeat association is not confined to origins; they are not origin
machinery and carry no implication for \emph{oriT} or \emph{oriV} biology.

These maps report model sensitivity, not biology. They do not measure
deletion phenotype, biochemical essentiality, binding affinity, physical
contact, or evolutionary conservation, and pattern intensity does not rank
functional importance. At pIGRK the outer inverted repeat IR1, which
published deletion analysis reports as dispensable for conjugative transfer,
carries a prominent pattern, while IR2, whose left arm is reported as
mandatory, does not carry a more prominent one
(\fref{fig:interpretation}a). At the NAH7 origin of transfer (CP096582.1,
[20737,20852)) the inverted repeat IR2 shows essentially no anti-diagonal
signal even though it is a perfect palindrome in the reading frame used ---
a frame that was itself selected to maximize arm self-complementarity, so
this negative is conditional on that choice
(\suppfref{fig:supp-orit-maps}c). We report it as an exploratory negative
rather than smoothing it over. Element recovery was not measured
systematically against a denominator, so no recovery rate is claimed.

One interval carries a pattern without a corresponding annotation. At pIGRK,
the 33 bp interval [1119,1152) shows an anti-diagonal pattern consistent with
self-complementarity but is not annotated in the source literature
(\fref{fig:interpretation}a, IR4). Its arms match at 9 of 10 positions and
the match degrades on extension, and its 13 bp loop,
\texttt{GAGCGATAGCGAA}, contains two of the four 5$'$-AGCGA-3$'$ motifs that
the source proposes as MobK recognition coordinates. This interval was
identified after inspecting the dependency map, by a stem-loop search over
the surrounding sequence rather than by a prespecified detector, and against
a shuffled null the arm-identity score is not significant (\emph{P} $\approx$
0.08). It may motivate a hairpin hypothesis, which we state as explicitly
unconfirmed: the categorical Jacobian does not measure physical folding, and
testing the hypothesis would require structure probing and mutagenesis
separating the stem from the AGCGA sequence itself. Any MobK-recognition role
is speculative. A comparable uncurated pattern occurs near position 30,480 in
the pHTbeta origin of transfer (\suppfref{fig:supp-orit-maps}a).

\edfref{fig:ed-umap-feature} and Supplementary
Figures~\ref{fig:supp-umap-detailed}--\ref{fig:supp-umap-gc}
% NOTE: the artifact reads "Supplementary Figs. S5-S7" -- one contiguous range
% over the three UMAP panels. The first of those was promoted to Extended Data
% (2026-09-01), so the range had to be split across the two bands.
show a two-dimensional UMAP
% TODO cite: UMAP.
of embeddings from the same frozen model, over 308,478 10 bp tiles from 40
origin-enriched Enterobacterales plasmids. The projection shows associations
with coarse functional annotation, with a finer annotation vocabulary, and
with GC content measured on the same tiles. Three limits apply. No separation
metric or classifier was computed, so we describe association rather than
class separation. The cohort was selected for curated \emph{oriT} and
\emph{oriV} hits and is not a representative sample of plasmids. And because
the projection varies strongly with GC content, sequence composition may
contribute to the apparent annotation structure. These embeddings come from
the final encoder layer, whereas the origin classifier described in
\secref{sec:origin} is fitted
% NOTE (2026-09-01): the artifact read "whereas the §3 classifier is fitted on
% blocks.60", and this line carried a TODO because \secref renders the heading
% TEXT -- giving "the Origin classification classifier". The §3 split made the
% collision worse, so this was recast on the pattern settled in the §2 split:
% the heading is now the object of "in", where a title reads naturally.
% sec:origin labels §3.1, which is where the classifier is defined.
on \texttt{blocks.60}, twelve residual blocks earlier; the two are different
representations of the same frozen model, and this projection is offered as
descriptive evidence, not as an account of classifier behaviour.

Together these selected examples show that pretraining on plasmid sequence
alone yields recurring model-derived dependencies associated with repeat
periodicity, self-complementarity, promoter organization and periodic
direct-repeat arrays, and that the same maps can flag candidate intervals,
such as the pIGRK [1119,1152) pattern, that carry no existing annotation and
can be tested experimentally.
